\documentclass[10pt,a4paper,twocolumn]{article}

\usepackage[utf8]{inputenc}
\usepackage[T1]{fontenc}
\usepackage{geometry}
\usepackage{parskip}
\usepackage{graphicx}
\usepackage{setspace}
\usepackage{amsmath}
\usepackage{hyperref}
\usepackage{subcaption}
\usepackage{float}

\geometry{margin=0.75in}
\setstretch{0.95}
\parskip=0.3em

\title{IPFS Network Analysis Report}
\author{}
\date{\today}

\begin{document}

\maketitle

\section{Introduction}

This report analyzes IPFS network telemetry generated by the Nebula crawler from March 14 to March 21. To account for geographical vantage point bias and the impact of local network middleboxes, the experiment utilized two distinct environments: (1) Consumer Edge: a MacBook Pro on a residential network in China, providing a "peer-level" view within a region known for complex NAT traversal and relay dependency; and (2) Cloud Infrastructure: a Japan-based VPS providing a stable, high-uptime perspective from a major Asian network hub.

A comparative analysis of the two crawl environments reveals a consistent global distribution ratio but a stark difference in absolute peer discovery. The local China-based crawler discovered 26,233 total multi-addresses—nearly double the 15,144 identified by the Japan server. Specifically, the local node identified 13,428 peers within China, while the remote server observed only 6,229, suggesting that regional network conditions and NAT traversal significantly impact node visibility.
While absolute peer discovery varies significantly between the two environments, the underlying metrics remain almost identical across both experiments. Consequently, to ensure depth and granularity, \textbf{all subsequent metrics and figures} in this report are derived from the local MacBook Pro dataset.
Insipired by Trautwein et al. (SIGCOMM 2022) paper, we analyze the IPFS network from the following aspects. Figure~\ref{fig:global} shows the geographical distribution of peer count per country. 
Figure~\ref{fig:global-peer-neighbour}  illustrates a highly concentrated distribution of peer connections, with the vast majority of IPFS nodes maintaining between 170 and 185 neighbors.
Figure~\ref{fig:global-new-found} tracks the fluctuating count of newly discovered peers after each crawl, showing a declining trend.
Figure~\ref{fig:multi-hash-prefix-asn} illustrates that the peer distribution is heavily skewed toward high-ranking networks; the top 10 ASes by CAIDA AS Rank account for 47.2\% of all peers, signaling a degree of centralization.
\mbox{} % Figures moved to last page to control placement.


\section{Peer ID Version}
Analyzing the ratio of v0 to v1 nodes identifies the network's technical maturity and cryptographic agility by distinguishing between legacy and modern identifiers. Based on crawler results, the network has overwhelmingly migrated to the newer CID/PeerID format: 93.11\% of nodes use the newer format, while only 6.88\% remain on the legacy format. However, v1 adoption is not uniform, and we can categorize regions by their ``modernity.'' This data (Table~\ref{tab:peerid_country}) provides a snapshot of the IPFS network's transition from \textbf{v0 (Qm)} to \textbf{v1 (12D3)} PeerIDs.
\begin{table}[ht]
    \centering
    \caption{Top 10 countries (total nodes $> 100$), sorted by v0 (Qm) share}
    \label{tab:peerid_country}
    \begin{tabular}{lrrrr}
        \hline
        Country & Qm & 12D3 & Total & Qm / Total \\
        \hline
        TW & 109 & 70  & 179  & 60.89\% \\
        HK & 112 & 88  & 200  & 56.00\% \\
        KR & 196 & 193 & 389  & 50.39\% \\
        UA & 24  & 98  & 122  & 19.67\% \\
        RU & 49  & 279 & 328  & 14.94\% \\
        SE & 15  & 87  & 102  & 14.71\% \\
        CH & 14  & 100 & 114  & 12.28\% \\
        US & 456 & 3883 & 4339 & 10.51\% \\
        PL & 12  & 105 & 117  & 10.26\% \\
        SG & 31  & 277 & 308  & 10.06\% \\
        \hline
    \end{tabular}
\end{table}


\section{Protocol Distribution}
To characterize the roles played by discovered peers, we analyzed the application protocols from the crawler. Across regions, the crawler most frequently observed \textbf{foundational IPFS/libp2p infrastructure}---basic connectivity and discovery (e.g., \texttt{/ipfs/ping}, \texttt{/ipfs/id}), routing and peer discovery (Kademlia DHT), and content exchange (Bitswap). 

Comparing the protocol distributions between China (CN) and the United States (US) reveals two distinct network profiles. In the United States (Figure~\ref{fig:protocol-dist-us}), the distribution is dominated by core IPFS/libp2p primitives and a large share of Bitswap-capable peers, consistent with broad participation in content exchange and well-connected public infrastructure. China (Figure~\ref{fig:protocol-dist-cn}) shows comparatively stronger representation of NAT traversal and relay-related capabilities, aligning with a more restrictive networking environment where peers rely on relay-assisted connectivity to remain reachable.
Notably, the presence of \texttt{/sbptp/1.0.0} in the China sample indicates that the crawler encountered peers advertising a protocol previously associated with IPFS-based botnet activity (often referred to as IPStorm). 
\mbox{} % Figures moved to last page to control placement.

\section{Agent Version Distribution}

The health of the ecosystem is not merely a matter of node count, but of the specific software versions powering those nodes. By examining regional distribution, we can identify clusters of potential vulnerability. Table\ref{tab:agent-coverage-country} shows that across countries, agent availability is uneven: some regions expose many reachable peers but only a small fraction return an agent string. Table~\ref{tab:agent-topshare-country} reveals a fragmentation of deployments characterized by a mix of legacy software and vendor-branded clusters. In Asia, particularly across China, South Korea, and Taiwan, the dominance of \texttt{go-ipfs/0.8.0} suggests pockets of outdated infrastructure that risk protocol incompatibility and exposure to older implementation bugs. In the United States, substantial infrastructure appears associated with \texttt{kubo/0.22.0/.../gala.games} and Docker-based deployments, implying that a small number of operator templates can shape regional risk: unpatched vulnerabilities or misconfigurations may be amplified when many nodes share the same build and operational pattern.
Notably, on the server side, the most dominant agent in TW is \texttt{storm}, accounting for 36.17\% of observed nodes, indicating possible IPStorm botnet activity.
\begin{table}[ht]
    \centering
    \caption{Top 5 agent response count by country}
    \label{tab:agent-coverage-country}
    \begin{tabular}{lrrr}
        \hline
        Country & Total & With agent & Without agent \\
        \hline
        CN & 13,428 & 624 & 12,804 \\
        US & 4,340 & 2,956 & 1,384 \\
        FR & 1,107 & 855 & 252 \\
        DE & 1,057 & 731 & 326 \\
        GB & 1,018 & 383 & 635 \\
        FI & 548 & 388 & 160 \\
        \hline
    \end{tabular}
\end{table}


\begin{table}[ht]
    \centering
    \caption{Top dominant agent percentage (of known agents) by country}
    \label{tab:agent-topshare-country}
    \begin{tabular}{llrr}
        \hline
        Country &  Agent & Share \\
        \hline
        CN & \texttt{go-ipfs/0.8.0/...} & 71.47\% \\
        AE & \texttt{kubo/0.37.0/docker} & 51.11\% \\
        MX & \texttt{go-ipfs/0.8.0/...} & 49.02\% \\
        AR & \texttt{kubo/0.32.1/} & 46.77\% \\
        KR & \texttt{go-ipfs/0.8.0/...} & 45.16\% \\
        TW & \texttt{go-ipfs/0.8.0/...} & 40.21\% \\
        HK & \texttt{go-ipfs/0.8.0/...} & 30.16\% \\
        IN & \texttt{kubo/0.32.1/} & 29.81\% \\
        TH & \texttt{kubo/0.32.1/} & 28.40\% \\
        ES & \texttt{kubo/0.37.0/.../docker} & 24.37\% \\
        GB & \texttt{kubo/0.39.0/} & 20.37\% \\
        US & \texttt{kubo/0.22.0/.../gala.games} & 20.03\% \\
        FR & \texttt{kubo/0.22.0/.../gala.games} & 20.00\% \\
        JP & \texttt{kubo/0.22.0/.../gala.games} & 20.00\% \\
        \hline
    \end{tabular}
\end{table}

\section{Node Availability}
\label{sec:available-nodes}

To assess node reliability despite intermittent monitoring outages, we employ a concept \textbf{Relative Availability (RA)} ratio. Unlike absolute uptime, which requires a continuous 24/7 observation window, RA measures a peer's presence relative only to the periods when the Nebula crawler was active. This ensures that the reliability scores reflect the node's actual network behavior rather than the monitoring station's connectivity gaps.
Figure~\ref{fig:peer-uptime-percentage} shows that most reachable peers have a high RA. 
Figure~\ref{fig:peer-uptime-country} shows that although CN has the highest peer count, its RA is much lower (11th in the ranking), whereas nodes in the US are more stable.
$$RA = \frac{\sum \text{Individual Peer Uptime}}{\sum \text{Total Crawler Active Sessions}}$$
Our analysis reveals a troubling \textbf{security-reliability issue} within the IPFS network, identified through Relative Availability. While the network is anchored by high-availability "Pillar" nodes, our data shows that a significant portion of this stable backbone is composed of compromised or high-risk entities. 
Specifically, the consistent detection of the \texttt{/sbptp/} protocol among our most reliable peers (Table~\ref{tab:protocol_counts_global}) suggests that poteltial malicious actors are maintaining a persistent, high-uptime presence to facilitate proxying and command-and-control activities. This risk is further compounded by the prevalence of outdated and vulnerable agents, such as \texttt{kubo/0.18.1} and \texttt{go-ipfs} (Table~\ref{tab:reliable-agent-share-rai-09}), which serve as stable but exploitable entry points. The fact that these "reliable" nodes are running malicious protocols or vulnerable software indicates that the network's most dependable participants are also its primary vectors for systemic instability and attack.
\begin{table}[H]
    \centering
    \small
    \caption{Most frequently observed protocols \(RA \ge 0.9\)}
    \label{tab:protocol_counts_global}
    \begin{tabular}{lr}
        \hline
        Protocol & Count \\
        \hline
        \texttt{/ipfs/ping/1.0.0} & 4841 \\
        \texttt{/ipfs/id/1.0.0} & 4837 \\
        \texttt{/ipfs/id/push/1.0.0} & 4834 \\
        \texttt{/ipfs/kad/1.0.0} & 4360 \\
        \texttt{/libp2p/autonat/1.0.0} & 4339 \\
        \texttt{/libp2p/circuit/relay/0.2.0/stop} & 4312 \\
        \texttt{/ipfs/bitswap/1.0.0} & 4294 \\
        \texttt{/ipfs/bitswap/1.1.0} & 4294 \\
        \texttt{/ipfs/bitswap} & 4291 \\
        \texttt{/x/} & 4280 \\
        \texttt{/ipfs/bitswap/1.2.0} & 4249 \\
        \texttt{/libp2p/dcutr} & 4012 \\
        \texttt{/ipfs/lan/kad/1.0.0} & 4000 \\
        \texttt{/libp2p/circuit/relay/0.2.0/hop} & 3681 \\
        \texttt{/libp2p/autonat/2/dial-back} & 2027 \\
        \texttt{/libp2p/autonat/2/dial-request} & 2027 \\
        \texttt{/libp2p/circuit/relay/0.1.0} & 1154 \\
        \texttt{/p2p/id/delta/1.0.0} & 1108 \\
        \texttt{/meshsub/1.0.0} & 344 \\
        \texttt{/meshsub/1.1.0} & 343 \\
        \texttt{/floodsub/1.0.0} & 342 \\
        \texttt{/sbptp/1.0.0} & 341 \\
        \texttt{/meshsub/1.2.0} & 221 \\
        \hline
    \end{tabular}
\end{table}
\begin{table}[ht]
  \centering
  \small
  \caption{Reliable agent share at \(RA \ge 0.9\)}
  \label{tab:reliable-agent-share-rai-09}
  \begin{tabular}{p{0.72\linewidth}r}
      \hline
      Agent & Percent \\
      \hline
      \texttt{kubo/0.22.0/3f884d3/gala.games} & 17.13\% \\
      \texttt{kubo/0.18.1/675f8bd/docker} & 10.06\% \\
      \texttt{kubo/0.37.0/6898472/docker} & 6.68\% \\
      \texttt{kubo/0.39.0/} & 6.28\% \\
      \texttt{go-ipfs/0.8.0/48f94e2} & 5.23\% \\
      \texttt{kubo/0.32.1/} & 3.58\% \\
      \texttt{kubo/0.39.0/2896aed/docker} & 3.32\% \\
      \texttt{kubo/0.37.0/6898472} & 1.96\% \\
      \texttt{kubo/0.40.1/39f8a65/docker} & 1.79\% \\
      \hline
  \end{tabular}
\end{table}

\clearpage
\onecolumn

% Single float: uniform 4×2 grid so panels share width/height and stay on one page.
\newlength{\finalpagefigH}
\setlength{\finalpagefigH}{0.152\textheight}
\captionsetup{font=footnotesize,skip=3pt}
\captionsetup[subfigure]{font=tiny,skip=0pt}
\setlength{\abovecaptionskip}{2pt}
\setlength{\belowcaptionskip}{0pt}

\begin{figure*}[t] 
  \centering
  \begin{subfigure}[t]{0.49\textwidth}
    \centering
    \includegraphics[width=\linewidth,height=\finalpagefigH,keepaspectratio]{pics/global_geographical.png}
    \caption{Multiaddresses per country (bucketed).}
    \label{fig:global}
  \end{subfigure}\hfill
  \begin{subfigure}[t]{0.49\textwidth}
    \centering
    \includegraphics[width=\linewidth,height=\finalpagefigH,keepaspectratio]{pics/global_peer_neighbour.png}
    \caption{Peer neighbour distribution.}
    \label{fig:global-peer-neighbour}
  \end{subfigure}

  \begin{subfigure}[t]{0.49\textwidth}
    \centering
    \includegraphics[width=\linewidth,height=\finalpagefigH,keepaspectratio]{pics/global_new_found.png}
    \caption{New peers found over time.}
    \label{fig:global-new-found}
  \end{subfigure}\hfill
  \begin{subfigure}[t]{0.49\textwidth}
    \centering
    \includegraphics[width=\linewidth,height=\finalpagefigH,keepaspectratio]{pics/multi_hash_prefix_asn.png}
    \caption{ASN distribution of peers.}
    \label{fig:multi-hash-prefix-asn}
  \end{subfigure}

  \begin{subfigure}[t]{0.49\textwidth}
    \centering
    \includegraphics[width=\linewidth,height=\finalpagefigH,keepaspectratio]{pics/protocol_distribution_US.png}
    \caption{Protocol distribution for the United States.}
    \label{fig:protocol-dist-us}
  \end{subfigure}\hfill
  \begin{subfigure}[t]{0.49\textwidth}
    \centering
    \includegraphics[width=\linewidth,height=\finalpagefigH,keepaspectratio]{pics/protocol_distribution_CN.png}
    \caption{Protocol distribution for China.}
    \label{fig:protocol-dist-cn}
  \end{subfigure}

  \begin{subfigure}[t]{0.49\textwidth}
    \centering
    \includegraphics[width=\linewidth,height=\finalpagefigH,keepaspectratio]{pics/peer_uptime_percentage.png}
    \caption{Distribution of peer RA values (bucketed).}
    \label{fig:peer-uptime-percentage}
  \end{subfigure}\hfill
  \begin{subfigure}[t]{0.49\textwidth}
    \centering
    \includegraphics[width=\linewidth,height=\finalpagefigH,keepaspectratio]{pics/peer_uptime_country.png}
    \caption{Reliable peers by country at different RA thresholds.}
    \label{fig:peer-uptime-country}
  \end{subfigure}

  \caption{%
  Global map and crawl metrics (a--d); regional protocol distributions (e--f); Relative Availability (RA) analysis (g--h). 
  \textbf{Note for (e--f):} Protocol IDs are truncated for clarity. Major mappings include: 
  \textit{bitswap} (v1.0.0--1.2.0), 
  \textit{relay} (circuit relay v0.1.0--0.2.0), 
  \textit{kad} (DHT Kademlia), 
  and \textit{autonat} (dial-back/request services).
    }
  \label{fig:global-metrics}
  \label{fig:protocol-distributions}
  \label{fig:peer-uptime-metrics}
\end{figure*}

\section{Future Analysis}
By correlating specific protocols (sbptp, sfst, and sbst) with reported Agent Versions, we identified a cluster of 1,154 compromised peers—comprising 973 \textbf{go-ipfs/0.8.0} and 181 \textbf{storm} nodes—establishing a definitive empirical footprint for the InterPlanetary Storm (IPStorm) botnet. This analysis reveals a significant geographic concentration within the Asia-Pacific region (specifically CN, TW, and KR), necessitating further investigation into node behavior to determine if these entities are centrally controlled for coordinated attacks. Moving forward, applying this Protocol-Agent Matrix to other vulnerable versions, such as \textbf{kubo/0.18.0} and \textbf{kubo/0.22.0}, will be critical in developing targeted mitigation strategies to preserve the security and high-performance standards of the global IPFS ecosystem.

\section{Data}
local data:
\href{https://bafybeico6jh7tqh4iavib73lv6mgk263vq4wmb727frskv34okdgwt2lc4.ipfs.dweb.link?filename=results_local.tar.gz}{\texttt{\scriptsize QmTepp4mfnTRAfTjW2ytV1VoSv9hE8tLDUDDpb3nby6CWr}}

server data:
\href{https://bafybeichuhco23no4oeaqc25kti6vgj4mi2nvv3wu5qffi4wm6om5uv5wq.ipfs.dweb.link?filename=results_server_server.tar}{\texttt{\scriptsize QmTAGj3FAJVMbpVbCuRXwJGX1r5wJLX6kDxvwDy8cGwVN3}}

\mbox{} % Figures moved to last page to control placement.

\end{document}
